\chapter{Types, Shapes, and the Shape Algebra}
\label{app:shapes}

\chapterorient{\Cref{ch:tensors} fixed the vocabulary of tensors and shapes, and
\cref{ch:shapes} taught the reading discipline: a shape is a type, named groups
assert congruence, scalar promotion lets a real number quietly join a complex
computation. Those chapters argued by example and picture. This appendix supplies
the \emph{formal} object underneath them. We write the shape language as a
grammar; we state named-group resolution as a single unification rule; we give the
type-and-shape judgment $\Gamma \vdash e : \tau$ as a small set of inference rules;
we collect the primitive shape contracts into one reference table; and we cast
scalar promotion as a subtyping rule \lstinline[style=l4style]|real <: complex|.
Nothing here is new physics or new notation---it is the precise version of the
rules the main text reads informally, the companion to \cref{app:calculus}'s
grammar and reduction rules and the home for the equational laws of
\cref{app:laws}. Read it as reference: the rule you reach for when a shape
contract or a typing question must be settled exactly rather than by eye.}

\section{What this appendix formalizes}
\label{sec:b-scope}

\Cref{ch:tensors,ch:shapes} carry the conceptual load; this appendix carries the
formal load behind them. The split is deliberate. The main text earns its rules on
worked vectors and matrices and draws each contract as a picture, so that a reader
learns to \emph{read} a shape the way they read a type. Here we write those same
rules as grammar productions, inference rules, and a reference table, so that a
reader can \emph{settle} a shape question exactly---decide whether two shapes are
congruent, whether a call type-checks, which group binds where---without appeal to
intuition.

Four objects organize the appendix, each the formal counterpart of a main-text
idea:

\begin{itemize}[leftmargin=1.4em, nosep]
\item \textbf{The shape language} (\cref{sec:b-shapelang})---a BNF grammar for the
shape expression $\sigma$, with its literals, axis variables, algebraic
combinations, the rank-wildcard \lstinline[style=l4style]|...|, and the named-group
binding and use forms. This is the formal version of \cref{ch:tensors}'s bracket
notation and \cref{ch:shapes}'s named groups.
\item \textbf{Named-group resolution} (\cref{sec:b-resolve})---the precise
unification rule by which a use \lstinline[style=l4style]|$S| is matched against a
binding \lstinline[style=l4style]|(S: ...)|, the formal version of
\cref{ch:shapes}'s rank-agnostic congruence.
\item \textbf{The type-and-shape judgment} (\cref{sec:b-judgment})---the inference
rules of $\Gamma \vdash e : \tau$, including the primitive-op shape-contract rule
that ties shapes into typing.
\item \textbf{The primitive shape contracts} (\cref{sec:b-contracts}) and
\textbf{scalar promotion} (\cref{sec:b-promotion})---a reference table of the
primitive signatures, and the subtyping rule
\lstinline[style=l4style]|real <: complex|.
\end{itemize}

\begin{keyidea}
\textbf{Shapes are types; congruence is unification.} A tensor's shape is a
type-level expression, checked at every function boundary exactly as a type is. A
named shape group \lstinline[style=l4style]|(S: ...)| is a unification variable
over a run of axes of \emph{unknown rank}; a use \lstinline[style=l4style]|$S|
resolves by \emph{unifying} the labelled shape against the binding, modulo a small
equational theory on axis extents (commutative, associative $+$ and $\times$). The
whole shape discipline of this book is therefore one rule applied everywhere:
\emph{unify the named groups of a signature against the argument shapes, modulo
that equational theory, rank first.}
\end{keyidea}

\section{The shape language}
\label{sec:b-shapelang}

A tensor type is \lstinline[style=l4style]|Tensor[s]| for a shape expression
$\sigma$. A shape expression is a tuple of \emph{items}, each item an axis
extent---a literal, a named axis, an algebraic combination---or one of the two
rank-flexible constructs: the wildcard and the named group. We give the grammar,
then read each production.

\subsection{Grammar}
\label{sec:b-grammar}

A shape $\sigma$ is a bracketed sequence of \emph{group items} $\gamma$; a group
item is either a single axis extent or one of the rank-flexible forms.

\begin{lstlisting}[style=l4style]
sigma ::= n                      -- literal natural (a pinned extent, e.g. 3)
        | s                      -- axis variable (a named axis: N, m, H, ...)
        | sigma + sigma          -- axis arithmetic
        | sigma - sigma
        | sigma * sigma
        | sigma ^ n
        | [ gamma_1, ..., gamma_k ]  -- a shape: a sequence of group items

gamma ::= sigma                  -- one axis (literal, variable, or algebraic)
        | ...                    -- rank wildcard: zero or more unconstrained axes
        | (S: gamma_1, ..., gamma_j)  -- named group BINDING: binds S to this run
        | $S                     -- named group USE: a back-reference to S
\end{lstlisting}

\noindent The grammar is taken verbatim from the source calculus's
\lstinline[style=l4style]|DimExpr| algebra; \cref{ch:tensors,ch:shapes} use exactly
its shapes. Reading the productions:

\begin{description}[leftmargin=1.6em, style=nextline, nosep]
\item[Literal $n$.] A \emph{pinned} extent: an axis whose length is fixed in the
signature, like the $3$ in \lstinline[style=l4style]|Tensor[H, W, 3]| (a 3-vector
per cell). A call supplying any other extent at that position is rejected.
\item[Axis variable $s$.] A \emph{named} axis: a single letter ($N$, $m$, $n$,
$H$, $W$) standing for a size not fixed in the signature but required to
\emph{match} wherever the same name recurs. An axis variable is \emph{one} axis,
hence rank-bearing---this is the load-bearing fact behind the
\lstinline[style=l4style]|Tensor[N]| pitfall of \cref{ch:shapes}.
\item[Algebraic combination.] An extent built from others by $+,-,\times,{}^\wedge$.
These arise from reshapes and factorizations: flattening a height axis $h$ and a
width axis $w$ into one produces the extent \lstinline[style=l4style]|h*w|. The
matcher reasons about these \emph{modulo} a small algebra (\cref{sec:b-dimexpr}).
\item[{Shape tuple $[\gamma_1,\dots,\gamma_k]$.}] The shape itself: an ordered
sequence of group items. The \emph{number} of axes the tuple denotes is the
tensor's \emph{rank}---but a single \lstinline[style=l4style]|...| or
\lstinline[style=l4style]|(S: ...)| item may itself stand for \emph{several} axes,
so the count of \emph{items} is not in general the rank.
\item[Rank wildcard \code{...}.] Matches \emph{zero or more}
axes of any extent. It is the device that makes a signature rank-agnostic: it
declines to commit to how many axes fill its position.
\item[Named group, binding \code{(S: ...)}.] Binds the name
$S$ to the run of axes its body matches. The binding occurs \emph{exactly once},
at $S$'s first appearance.
\item[{Named group, use \code{\$S}.}] A \emph{back-reference}: it
asserts that the run of axes at this position is \emph{congruent} to the run bound
to $S$. The \lstinline[style=l4style]|$| sigil is what visually distinguishes a use
from a fresh binding or a single-axis variable.
\end{description}

\Cref{fig:b-shapealgebra} draws the shape algebra as labelled axis boxes: a row of
boxes for a tensor's axes, a spanning bracket for a named group's binding, and the
\lstinline[style=l4style]|$S| back-arrow for a use---the formal-side companion to
the hero figure of \cref{ch:shapes}.

\begin{figure}[t]
\centering
\begin{tikzpicture}[scale=1.0]
  % LEGEND-style: literal, variable, group, wildcard
  % Row of axis boxes with mixed item kinds: [ (S: a, ...), k ]
  \begin{scope}[shift={(0,1.3)}]
    % the named-group run: a leading pinned 'a' then a wildcard tail
    \draw[draw=simBlue, fill=simBgBlue, thick] (0,-0.4) rectangle (0.8,0.4);
    \node[font=\small, simInk] at (0.4,0) {$a$};
    \draw[draw=simBlue, fill=simBgBlue, thick] (1.0,-0.4) rectangle (1.8,0.4);
    \node[font=\footnotesize, simGray] at (1.4,0) {$\cdot$};
    \node[font=\scriptsize, simGray] at (2.15,0) {$\cdots$};
    \draw[draw=simBlue, fill=simBgBlue, thick] (2.5,-0.4) rectangle (3.3,0.4);
    \node[font=\footnotesize, simGray] at (2.9,0) {$\cdot$};
    % the free trailing axis k
    \draw[draw=simRust, fill=simBgRust, thick] (3.9,-0.4) rectangle (4.7,0.4);
    \node[font=\small, simInk] at (4.3,0) {$k$};
    % bracket binding the group S over the leading run
    \draw[decorate, decoration={brace, amplitude=4pt}, simBlue]
      (-0.02,0.55) -- (3.32,0.55)
      node[midway, above=4pt, font=\scriptsize\bfseries, simBlue]
      {binds \texttt{S} (a leading run: pinned \texttt{a}, then \texttt{...})};
    \node[font=\scriptsize\bfseries, simRust, above=4pt] at (4.3,0.55) {free \texttt{k}};
    \node[font=\scriptsize\ttfamily, simInk] at (-2.0,0)
      {\lstinline[style=l4style]|Tensor[(S: a, ...), k]|};
  \end{scope}
  % a second operand: $S then a free m  -- the USE row
  \begin{scope}[shift={(0,-0.6)}]
    \draw[draw=simTeal, fill=simBgGreen, thick] (0,-0.4) rectangle (0.8,0.4);
    \node[font=\small, simInk] at (0.4,0) {$a$};
    \draw[draw=simTeal, fill=simBgGreen, thick] (1.0,-0.4) rectangle (1.8,0.4);
    \node[font=\footnotesize, simGray] at (1.4,0) {$\cdot$};
    \node[font=\scriptsize, simGray] at (2.15,0) {$\cdots$};
    \draw[draw=simTeal, fill=simBgGreen, thick] (2.5,-0.4) rectangle (3.3,0.4);
    \node[font=\footnotesize, simGray] at (2.9,0) {$\cdot$};
    \draw[draw=simGold, fill=simBgGold, thick] (3.9,-0.4) rectangle (4.7,0.4);
    \node[font=\small, simInk] at (4.3,0) {$m$};
    \draw[decorate, decoration={brace, amplitude=4pt, mirror}, simTeal]
      (-0.02,-0.55) -- (3.32,-0.55)
      node[midway, below=4pt, font=\scriptsize\bfseries, simTeal]
      {uses \texttt{\$S} (must be congruent)};
    \node[font=\scriptsize\bfseries, simGold, below=4pt] at (4.3,-0.55) {free \texttt{m}};
    \node[font=\scriptsize\ttfamily, simInk] at (-2.0,0)
      {\lstinline[style=l4style]|Tensor[$S, m]|};
  \end{scope}
  % the congruence back-arrow
  \draw[-{Stealth[length=2.6mm]}, simRust, thick]
    (3.55,-0.85) .. controls (4.5,-0.1) and (4.5,0.7) .. (3.55,1.45)
    node[midway, right=1pt, font=\scriptsize\itshape, simRust, align=left]
    {congruent\\(same run)};
\end{tikzpicture}
\caption[The shape algebra: items, groups, and the congruence back-arrow]{The
shape language as labelled axis boxes. \emph{Top (blue/rust):} the binding shape
\lstinline[style=l4style]|Tensor[(S: a, ...), k]|---a named group \texttt{S}
spanning a leading run (a pinned axis \texttt{a} followed by a wildcard
\texttt{...} of any number of further axes), then a single \emph{free} trailing
axis \texttt{k} outside the group. \emph{Bottom (teal/gold):} a second operand
\lstinline[style=l4style]|Tensor[$S, m]|, whose leading run is a \emph{use}
\lstinline[style=l4style]|$S| and whose own trailing axis \texttt{m} is free and
independent of \texttt{k}. The rust back-arrow is the congruence assertion the
\lstinline[style=l4style]|$S| carries: \emph{my leading run is the same run that
bound \texttt{S}}. Group items may be literals, variables, algebraic extents,
wildcards, or nested groups; a group splits a shape into a named congruent part
and a free part.}
\label{fig:b-shapealgebra}
\end{figure}

\subsection{The DimExpr equational theory}
\label{sec:b-dimexpr}

Two shapes that differ on the page may denote the same run of extents. A reshape
that merges axes $h$ and $w$ yields the extent \lstinline[style=l4style]|h*w|; an
operation consuming the flattened tensor must recognize
\lstinline[style=l4style]|h*w| and \lstinline[style=l4style]|w*h| as the same
length. Shape matching is therefore \emph{not} raw character comparison; it is
comparison modulo a small equational theory on extent expressions---the
\lstinline[style=l4style]|DimExpr| theory.

\begin{definition}[DimExpr equational theory]
\label{def:b-dimexpr}
Two axis extents are \emph{equal} when they are equal as elements of the free
commutative semiring on the axis variables, with $+$ and $\times$ the semiring
operations: that is, modulo the axioms

\begin{center}
\begin{tabular}{@{}ll@{}}
\toprule
\emph{commutativity} & $a + b = b + a$,\qquad $a \times b = b \times a$ \\
\emph{associativity} & $(a+b)+c = a+(b+c)$,\qquad $(a\,b)\,c = a\,(b\,c)$ \\
\emph{distributivity} & $a\,(b+c) = a\,b + a\,c$ \\
\emph{identities} & $a + 0 = a$,\qquad $a \times 1 = a$,\qquad $a \times 0 = 0$ \\
\emph{powers} & $a^{m}\,a^{n} = a^{m+n}$,\qquad $a^{1}=a$ \\
\bottomrule
\end{tabular}
\end{center}

\noindent Subtraction $a - b$ is admitted as the formal inverse of addition
(extents arising from a reshape are assumed nonnegative; the theory does not
attempt to prove nonnegativity). Two shapes $[\gamma_1,\dots,\gamma_k]$ and
$[\delta_1,\dots,\delta_k]$ of \emph{equal rank} are equal when $\gamma_i = \delta_i$
in this theory for every $i$.
\end{definition}

In practice the matcher normalizes each extent expression to a canonical
sum-of-monomials form and compares; \lstinline[style=l4style]|h*w| and
\lstinline[style=l4style]|w*h| both normalize to the monomial $hw$, so they match,
while \lstinline[style=l4style]|h*w| and \lstinline[style=l4style]|h+w| do not.
The theory is deliberately weak---it decides the \emph{algebraic} coincidences a
reshape genuinely produces, and nothing subtler.

\begin{notationbox}
The equational theory is a fact the matcher \emph{verifies}; congruence is a
contract you \emph{assert}. When two operands must share a shape, say so with a
group (\lstinline[style=l4style]|(S: ...)| / \lstinline[style=l4style]|$S|) and let
the matcher enforce it; reserve the extent algebra
(\lstinline[style=l4style]|h*w| and friends) for the positions where a genuine
reshape \emph{produced} a composite extent. Writing two operands as
\lstinline[style=l4style]|Tensor[h*w]| and hoping the algebra makes them agree is
backwards: the algebra equates \emph{extents at matching positions}, it does not
discover that two whole shapes were meant to be congruent.
\end{notationbox}

\section{Named-group resolution}
\label{sec:b-resolve}

We can now state, precisely, what \cref{ch:shapes} drew as a back-arrow: how a use
\lstinline[style=l4style]|$S| is resolved against a binding
\lstinline[style=l4style]|(S: ...)|. The resolution is a \emph{unification} of
axis runs, with one wrinkle---a named group spans a \emph{run} of axes of unknown
length, so the unifier must split a concrete shape into the part the group claims
and the part left free.

\subsection{The resolution rule}
\label{sec:b-resolve-rule}

A signature is checked left to right against the argument shapes. The first
position labelled by a group binds it; later uses are tested.

\begin{definition}[Named-group resolution]
\label{def:b-resolve}
Let a signature contain a binding \lstinline[style=l4style]|(S: p)| (where the
\emph{pattern} $p$ is a sequence of group items, possibly containing a wildcard
\lstinline[style=l4style]|...|) and one or more uses \lstinline[style=l4style]|$S|.
Resolution proceeds in two steps.

\begin{description}[leftmargin=1.6em, style=nextline, nosep]
\item[Bind.] Match the pattern $p$ against the corresponding run of the
\emph{first} argument shape that labels $S$. If $p$ is a pure wildcard
\lstinline[style=l4style]|(S: ...)|, the run is the \emph{entire} shape at that
position; if $p$ pins leading or trailing axes (e.g.
\lstinline[style=l4style]|(S: a, ...)|), those pinned axes must match by the
\lstinline[style=l4style]|DimExpr| theory (\cref{def:b-dimexpr}) and the wildcard
absorbs the remainder. The matched run---a concrete sequence of extents
$\hat{S} = [\hat{e}_1,\dots,\hat{e}_r]$ of definite rank $r$---becomes the
\emph{value} of $S$.
\item[Use.] At each later \lstinline[style=l4style]|$S|, require the run of axes at
that position to be \emph{congruent} to $\hat{S}$: \textbf{(i)} the same rank $r$,
and \textbf{(ii)} extent $\hat{e}_i$ at axis $i$ for every $i$, equality decided by
the \lstinline[style=l4style]|DimExpr| theory. Rank is checked \emph{first}: a rank
mismatch fails immediately, before any extent is compared.
\end{description}

\noindent A group with both a pinned part and a free part splits each labelled
shape into ``the run $S$ claims'' and ``the free axes outside it''; only the
claimed run participates in congruence, while the free axes (\texttt{k},
\texttt{m} in \cref{fig:b-shapealgebra}) are unconstrained across operands.
\end{definition}

Read as a unification rule, the binding/use machinery is one inference:

\[
\resizebox{\linewidth}{!}{$\displaystyle
\dfrac{
  \begin{array}{c}
  \text{$\hat{S}$ is the run matched by \lstinline[style=l4style]|(S: p)| in the binding position}
  \end{array}
  \qquad
  \text{rank}(\tau) = \text{rank}(\hat{S})
  \qquad
  \tau \;=_{\text{DimExpr}}\; \hat{S}
}{
  \text{a shape $\tau$ at a use position \lstinline[style=l4style]|$S| resolves (is congruent)}
}
\quad(\textsc{congr})
$}
\]

\noindent and the dual---a rank mismatch---is the rejection rule:

\[
\dfrac{\text{rank}(\tau) \neq \text{rank}(\hat{S})}
      {\text{$\tau$ at \lstinline[style=l4style]|$S| is rejected (rank mismatch, no extent compared)}}
\quad(\textsc{congr-fail})
\]

This is exactly why \lstinline[style=l4style]|Tensor[N]| (\cref{ch:shapes}) can
never match a rank-$2$ argument: a single axis variable $N$ is a run of rank $1$,
so \textsc{congr-fail} fires on the rank check before $N$'s extent is ever
considered. \Cref{fig:b-unify} draws the two-step resolution as a matched
binding-against-use, the formal version of the boundary gate of
\cref{ch:shapes}.

\begin{figure}[t]
\centering
\begin{tikzpicture}[node distance=8mm and 12mm]
  % the binding shape on the left
  \node[data, text width=26mm] (bind)
    {\textbf{binding}\\\lstinline[style=l4style]|(S: ...)|\\matched run\\$\hat{S}=[3,4]$};
  % the unify gate
  \node[stage, right=16mm of bind, minimum height=24mm, text width=20mm] (gate)
    {\textbf{unify}\\$\$S \stackrel{?}{=} \hat{S}$};
  % two use candidates
  \node[data, right=16mm of gate, yshift=11mm, text width=30mm] (ok)
    {use \lstinline[style=l4style]|$S| $=[3,4]$\\rank $2$, extents agree};
  \node[data, right=16mm of gate, yshift=-11mm, text width=30mm] (bad)
    {use \lstinline[style=l4style]|$S| $=[12]$\\rank $1\neq 2$};
  \draw[flow] (bind) -- node[above, font=\scriptsize, simGray]{bind $\hat{S}$} (gate);
  \draw[-{Stealth[length=2.4mm]}, simGreen, thick] (gate.east|-ok) -- (ok)
    node[midway, above, font=\scriptsize, simGreen]{congruent};
  \draw[-{Stealth[length=2.4mm]}, simRust, thick, densely dashed] (gate.east|-bad) -- (bad)
    node[midway, below, font=\scriptsize, simRust]{reject};
  % annotation: rank first
  \node[font=\scriptsize\itshape, simGray, below=7mm of gate, text width=64mm,
        align=center]
    {\textbf{step 1} bind $\hat{S}$ from the first labelled run;\quad
     \textbf{step 2} at each use, check rank \emph{first}, then extents modulo
     \lstinline[style=l4style]|DimExpr|. The lower candidate fails on rank
     ($1\neq 2$) before any extent is compared---even though $3\cdot4=12$.};
\end{tikzpicture}
\caption[Named-group resolution as unification]{Resolving a use
\lstinline[style=l4style]|$S| against a binding \lstinline[style=l4style]|(S: ...)|.
The binding's first labelled run fixes $\hat{S}$ (here $[3,4]$, rank $2$). Each
later use is \emph{unified} against $\hat{S}$: rank is checked first (a mismatch is
rejected on the spot, lower path), then extents axis-by-axis modulo the
\lstinline[style=l4style]|DimExpr| equational theory (\cref{def:b-dimexpr}, upper
path). The flattened candidate $[12]$ is rejected on \emph{rank} even though its
single extent equals the product $3\cdot4$ of the bound run: congruence demands the
same number of axes, not merely the same total size. This is the precise content of
\cref{ch:shapes}'s back-arrow and boundary gate.}
\label{fig:b-unify}
\end{figure}

\subsection{Operator shapes: range-first}
\label{sec:b-operators}

An operator whose domain and range shapes differ names \emph{two} groups, a range
group $R$ and a domain group $D$, written \emph{range-first} (matching the
$m \times n$ matrix convention, rows then columns):

\begin{lstlisting}[style=l4style]
LinOp[(R: ...), (D: ...)]                                   -- general operator
apply_linop :: LinOp[(R: ...), (D: ...)] -> Tensor[$D] -> Tensor[$R]
\end{lstlisting}

\noindent The operand is a use of the \emph{domain} group
(\lstinline[style=l4style]|$D|); the result is a use of the \emph{range} group
(\lstinline[style=l4style]|$R|). Resolution is \cref{def:b-resolve} applied to two
groups at once: the operator value binds both $R$ and $D$, the operand is checked
against $\hat{D}$, the result is congruent to $\hat{R}$. An \emph{endomorphism}
(domain $\equiv$ range) binds one group and uses it for both:

\begin{lstlisting}[style=l4style]
LinOp[(S: ...), $S]                                        -- square / endomorphic
apply_linop :: LinOp[(S: ...), $S] -> Tensor[$S] -> Tensor[$S]
\end{lstlisting}

\noindent which is the signature a Krylov solver wants: feed it a vector of shape
$S$, get back a vector of shape $S$ you can feed it again
(\cref{ch:matrixfree}). At the lowest layers, where Palace's operators act on
genuine flat dof-vectors of a definite length, the concrete rank-$1$ spelling
\lstinline[style=l4style]|LinearOperator[M, N]| over
\lstinline[style=l4style]|Tensor[N]| is faithful and is kept; the group form is the
rank-agnostic calculus rendering used at the higher layers.

\subsection{Congruence groups versus named axes of fixed meaning}
\label{sec:b-fixedaxes}

The resolution rule applies only to \emph{congruence groups}---the
\lstinline[style=l4style]|(S: ...)| / \lstinline[style=l4style]|$S| variables. A
shape bracket holds a second, entirely different kind of name, and the two must not
be confused. A \emph{named axis of fixed meaning} is a single axis whose letter
denotes one specific, fixed quantity---it is \emph{not} a unification variable and
\emph{not} a back-reference. The book's recurring family is the element-local
vocabulary of the matrix-free kernel (\cref{ch:matrixfree}):

\begin{description}[leftmargin=1.6em, style=nextline, nosep]
\item[\texttt{E}] element count---how many mesh elements the operator runs over.
\item[\texttt{L}] local dofs per element, $(p+1)^d$ for a degree-$p$ element.
\item[\texttt{P}] quadrature points per element.
\item[\texttt{C}] value (or derivative) components carried at a quadrature point.
\item[\texttt{G}] geometry-factor components stored per quadrature point.
\end{description}

\noindent with the three rank-structured shapes
\lstinline[style=l4style]|Tensor[(E, L)]|, \lstinline[style=l4style]|Tensor[(E, P, C)]|,
and \lstinline[style=l4style]|Tensor[(E, P, G)]|. Each letter has a definite role
fixed by the finite-element construction; none is rank-agnostic, none is a
back-reference. A fixed-meaning axis answers \emph{what is this axis}; a congruence
group answers \emph{does this shape match that one}. The two are governed by
different rules: a fixed-meaning axis is matched like an ordinary axis variable
(its letter must agree wherever it recurs, at rank $1$), whereas a congruence group
is resolved by \cref{def:b-resolve} over a run of unknown rank. The flat global
dof-axis $N$ is a third species again---a genuine rank-$1$
\lstinline[style=l4style]|Tensor[N]|, neither fixed-meaning element-local nor a
congruence group---and the gather/scatter pair $G,G^{\mathsf T}$ of
\cref{ch:matrixfree} is precisely the boundary between the flat-$N$ world and the
fixed-meaning \lstinline[style=l4style]|(E, L)| world.

\section{The type-and-shape judgment}
\label{sec:b-judgment}

Shapes do not float free; they ride on \emph{types}, and a tensor type
\lstinline[style=l4style]|Tensor[s]| is a type former carrying a shape. The
calculus's typing judgment is the standard typed-$\lambda$ form,
\[
  \Gamma \vdash e : \tau,
\]
read ``under context $\Gamma$ (a finite map from variables to types), the
expression $e$ has type $\tau$.'' The type $\tau$ ranges over the ground scalar
types, tensor types, records, tuples, functions, the simulator monad
\lstinline[style=l4style]|Sim t|, and operator types
\lstinline[style=l4style]|Op[t_in -> t_out]| (the full grammar is
\cref{app:calculus}). The shape discipline of the previous sections enters typing
at exactly one place: the rule for primitive-operator application, where the op's
declared shape contract must match the argument shapes modulo the equational theory.

\subsection{The core rules}
\label{sec:b-coretyping}

The structural rules are the textbook ones. We give variable, application,
abstraction, record formation and spread, and the two monadic rules; the rest are
mechanical.

\begin{description}[leftmargin=1.4em, style=nextline]
\item[Variable.] A variable has the type the context assigns it.
\[
  \dfrac{x : \tau \,\in\, \Gamma}{\Gamma \vdash x : \tau}\quad(\textsc{var})
\]

\item[Application and abstraction.] Function application demands the argument's
type meet the function's domain; abstraction discharges a hypothesis.
\[
  \dfrac{\Gamma \vdash e_1 : \tau_1 \to \tau_2 \qquad \Gamma \vdash e_2 : \tau_1}
        {\Gamma \vdash e_1\,e_2 : \tau_2}\quad(\textsc{app})
  \qquad\qquad
  \dfrac{\Gamma,\, x : \tau_1 \vdash e : \tau_2}
        {\Gamma \vdash \lambda(x{:}\tau_1).\,e \;:\; \tau_1 \to \tau_2}\quad(\textsc{abs})
\]

\item[Records and spread.] A record literal is typed field-by-field; a spread
(functional update) preserves the record type while replacing one field's value
with one of the same type.
\[
  \dfrac{\Gamma \vdash e_i : \tau_i \ \ (\text{each } i)}
        {\Gamma \vdash \{\, l_1{:}\,e_1,\dots,l_n{:}\,e_n \,\} : \{\, l_1{:}\,\tau_1,\dots,l_n{:}\,\tau_n \,\}}
  \quad(\textsc{record})
\]
\[
  \dfrac{\Gamma \vdash e : \{\, l_1{:}\,\tau_1,\dots,l_n{:}\,\tau_n \,\}
         \qquad \Gamma \vdash w : \tau_k}
        {\Gamma \vdash \{\, \dots e,\ l_k{:}\,w \,\} : \{\, l_1{:}\,\tau_1,\dots,l_n{:}\,\tau_n \,\}}
  \quad(\textsc{spread})
\]

\item[Monad: return and bind.] \lstinline[style=l4style]|return| injects a pure
value into the simulator monad; bind sequences two \lstinline[style=l4style]|Sim|
actions, threading the state (the ownership discipline of \cref{ch:monad}).
\[
  \dfrac{\Gamma \vdash e : \tau}{\Gamma \vdash \texttt{return}\ e : \texttt{Sim}\ \tau}
  \quad(\textsc{return})
  \qquad\quad
  \dfrac{\Gamma \vdash m : \texttt{Sim}\ \alpha \qquad \Gamma,\, x{:}\,\alpha \vdash k : \texttt{Sim}\ \beta}
        {\Gamma \vdash m \mathbin{>\!\!>\!\!=} \lambda x.\,k \;:\; \texttt{Sim}\ \beta}
  \quad(\textsc{bind})
\]
\end{description}

\noindent These rules are shape-agnostic: a shape appears only \emph{inside} a
tensor type and is carried along opaquely by \textsc{var}, \textsc{app}, and the
rest. Where shapes acquire force is the next rule.

\subsection{The primitive shape-contract rule}
\label{sec:b-primrule}

Each primitive operator $\mathit{op}$ carries a declared signature---a sequence of
argument tensor (or scalar) types and a result type, with shapes drawn from
\cref{sec:b-shapelang}, possibly mentioning groups. Applying $\mathit{op}$ to
arguments type-checks when each argument's shape matches the contract's
corresponding shape \emph{modulo the equational theory and the group resolution of
\cref{def:b-resolve}}, after which the result wears the contract's result shape with
the resolved groups substituted in.

Write the contract of $\mathit{op}$ as
\lstinline[style=l4style]|op :: Tensor[s_1] -> ... -> Tensor[s_k] -> Tensor[s_r]|,
possibly with scalar arguments interspersed and with groups shared among the
$\sigma$'s. Let $\theta$ be the substitution that resolves every group
(\cref{def:b-resolve}) against the argument shapes. The rule is:

\[
\dfrac{
  \begin{array}{c}
  \mathit{op} :: \texttt{Tensor[}\sigma_1\texttt{]} \to \cdots \to \texttt{Tensor[}\sigma_k\texttt{]} \to \texttt{Tensor[}\sigma_r\texttt{]}
  \qquad
  \Gamma \vdash e_i : \texttt{Tensor[}\tau_i\texttt{]} \ \ (\text{each } i) \\[2pt]
  \tau_i \;=_{\text{DimExpr}}\; \theta(\sigma_i) \ \ (\text{each } i)
  \qquad
  \theta \text{ resolves every group of the contract (\cref{def:b-resolve})}
  \end{array}
}{
  \Gamma \vdash \mathit{op}(e_1,\dots,e_k) : \texttt{Tensor[}\theta(\sigma_r)\texttt{]}
}\quad(\textsc{prim})
\]

\noindent In words: the op's declared shape contract must match the argument shapes
modulo the equational theory (with groups resolved consistently), and the result
shape is the contract's result with the resolved group values plugged in. This is
the single rule that turns ``a shape is a type'' into a checkable obligation---and
it is the formal content of every shape-matching picture in
\cref{ch:tensors,ch:shapes}. A mismatch in any $\tau_i =_{\text{DimExpr}}
\theta(\sigma_i)$, or a failed group resolution, is a type error caught at the
boundary, before any arithmetic runs.

\begin{example}[\code{matvec} under \textsc{prim}]
\label{ex:b-matvec}
The contract is \lstinline[style=l4style]|matvec :: Tensor[m, n] -> Tensor[n] -> Tensor[m]|.
Apply it to a matrix of shape $[3,4]$ and a vector of shape $[4]$. Resolution
binds $m \mapsto 3$, $n \mapsto 4$ from the first argument; the second argument's
shape $[4]$ must equal $\theta([n]) = [4]$---it does. By \textsc{prim} the result
type is $\theta([m]) = [3]$, a length-$3$ vector. Hand it a length-$3$ vector
instead and the premise $[3] =_{\text{DimExpr}} [4]$ fails: a type error at the
boundary, exactly the shape-mismatch picture of \cref{ch:tensors}.
\end{example}

\Cref{fig:b-deriv} shows \textsc{prim} composing with the structural rules in a
small derivation tree: typing the composite expression
\lstinline[style=l4style]|axpy a x (matvec K x)|, where the inner
\lstinline[style=l4style]|matvec| produces the second \lstinline[style=l4style]|axpy|
operand.

\begin{figure}[t]
\centering
\begin{tikzpicture}[scale=1.0,
   every node/.style={font=\footnotesize}]
  % leaves
  \node (g1) at (-5.4,2.4) {$\Gamma \vdash \mathsf{K} : \texttt{Tensor[3,3]}$};
  \node (g2) at (-2.6,2.4) {$\Gamma \vdash \vx : \texttt{Tensor[3]}$};
  \node (g3) at (-2.6,0.9) {$\Gamma \vdash \vx : \texttt{Tensor[3]}$};
  \node (ga) at (-5.0,0.9) {$\Gamma \vdash a : \texttt{Scalar}$};
  % matvec node (PRIM)
  \node (mv) at (-3.9,1.55)
    {$\Gamma \vdash \texttt{matvec}\ \mathsf{K}\ \vx : \texttt{Tensor[3]}$};
  \draw[simGray] (-5.4,2.2) -- (-2.0,2.2);
  \node[font=\scriptsize, simRust] at (-1.5,2.2) {(\textsc{prim})};
  % axpy node (PRIM), top conclusion
  \node (ax) at (-3.4,-0.3)
    {$\Gamma \vdash \texttt{axpy}\ a\ \vx\ (\texttt{matvec}\ \mathsf{K}\ \vx) : \texttt{Tensor[3]}$};
  \draw[simGray] (-5.4,0.65) -- (-1.6,0.65);
  \node[font=\scriptsize, simRust] at (-1.1,0.65) {(\textsc{prim})};
  % connectors (informal tree edges)
  \draw[simGray, -] (mv.south) -- (-3.9,0.7);
  \node[font=\scriptsize\itshape, simGray, align=center] at (1.9,1.55)
    {\textbf{inner} \texttt{matvec}\\binds $m,n\mapsto 3$;\\result \texttt{Tensor[3]}};
  \node[font=\scriptsize\itshape, simGray, align=center] at (1.9,-0.3)
    {\textbf{outer} \texttt{axpy}\\binds $\$S\mapsto[3]$ from $\vx$;\\
     the \texttt{matvec} result, also\\\texttt{[3]}, is congruent};
  \draw[refedge, simBlue] (0.3,1.55) -- (mv.east);
  \draw[refedge, simBlue] (0.3,-0.3) -- (ax.east);
\end{tikzpicture}
\caption[A typing derivation for a composite expression]{A derivation of
$\Gamma \vdash \texttt{axpy}\ a\ \vx\ (\texttt{matvec}\ \mathsf{K}\ \vx) :
\texttt{Tensor[3]}$, showing \textsc{prim} (the primitive shape-contract rule)
applied twice and composed through the result of one op feeding the next. The
\emph{inner} \lstinline[style=l4style]|matvec| (top) takes the $[3,3]$ operator
$\mathsf{K}$ and the $[3]$ vector $\vx$, binds $m,n\mapsto 3$, and yields a $[3]$
result by \textsc{prim} (\cref{ex:b-matvec}). That $[3]$ result becomes the third
operand of the \emph{outer} \lstinline[style=l4style]|axpy| (bottom), whose group
\lstinline[style=l4style]|$S| is bound to $[3]$ by $\vx$; the
\lstinline[style=l4style]|matvec| result, also $[3]$, is congruent, so the whole
expression types at \lstinline[style=l4style]|Tensor[3]|. Each horizontal bar is one
rule instance; the tree is read bottom-up as ``to conclude the root, discharge these
premises.''}
\label{fig:b-deriv}
\end{figure}

\section{The primitive shape contracts: a reference}
\label{sec:b-contracts}

We collect the primitive operators' shape contracts in one place. Each is a pure
function (\cref{ch:operators} gives their algebra and laws; \cref{app:laws}
collects the equational laws). The table is the formal companion to the
shape-matching pictures of \cref{ch:tensors}: read each row as a signature plus the
side-condition its groups and axis names impose under \textsc{prim}.

\begin{figure}[t]
\centering
\small
\renewcommand{\arraystretch}{1.35}
\resizebox{\linewidth}{!}{%
\begin{tabular}{@{}>{\ttfamily}l l l@{}}
\toprule
\normalfont\textbf{op} & \textbf{signature} & \textbf{shape side-condition} \\
\midrule
axpy &
  \lstinline[style=l4style]|Scalar -> Tensor[$S] -> Tensor[$S] -> Tensor[$S]| &
  all three congruent ($S$) \\
dot &
  \lstinline[style=l4style]|Tensor[$S] -> Tensor[$S] -> Scalar| &
  both congruent; collapses \emph{all} axes \\
matvec &
  \lstinline[style=l4style]|Tensor[m, n] -> Tensor[n] -> Tensor[m]| &
  inner $n$ must match, cancels; $m$ survives \\
outer &
  \lstinline[style=l4style]|Tensor[m] -> Tensor[n] -> Tensor[m, n]| &
  no axis matches; rank \emph{raised} ($1{+}1$) \\
reduce &
  \lstinline[style=l4style]|Tensor[..., k] -> Tensor[...]| &
  labelled axis $k$ collapsed; rest survives \\
unfold &
  \lstinline[style=l4style]!Tensor[$S] -> (axis,size,stride) -> Tensor[..., size]! &
  windows one axis into a new trailing \texttt{size} \\
broadcast &
  \lstinline[style=l4style]!(target: Shape) -> Tensor[$Sa] -> Tensor[target]! &
  \texttt{target} a broadcast extension of $S_a$ \\
conv &
  \lstinline[style=l4style]|Tensor[N,Ci,Si] -> Tensor[Co,Ci,Sk] -> Tensor[N,Co,So]| &
  channels $C_i$ match; $S_o$ from $S_i,S_k$ \\
\bottomrule
\end{tabular}%
}
\caption[Primitive shape contracts]{The primitive operators and their shape
contracts, the formal reference behind \cref{ch:tensors}'s picture-per-op
treatment. \emph{Same-shape} ops (\lstinline[style=l4style]|axpy|,
\lstinline[style=l4style]|dot|) demand their operands be congruent over a group
$S$; \emph{contraction} (\lstinline[style=l4style]|matvec|) and
\emph{rank-raising} (\lstinline[style=l4style]|outer|) ops state the input/output
axis relation explicitly; \emph{rank-moving} ops
(\lstinline[style=l4style]|reduce| drops an axis,
\lstinline[style=l4style]|broadcast| / \lstinline[style=l4style]|unfold| add one);
and \lstinline[style=l4style]|conv| is the structured kernel contract (matching
input channels $C_i$, output spatial extent $S_o$ a function of input and kernel
extents $S_i$, $S_k$). Each side-condition is what the \textsc{prim} rule
(\cref{sec:b-primrule}) checks at the boundary. In running text we abbreviate the
broadcast/unfold subscript names; the source signatures use
\lstinline[style=l4style]!Tensor[$S_a]! and the explicit \texttt{target\_shape}
tuple.}
\label{fig:b-contracts}
\end{figure}

A few rows reward a second look, each a representative \emph{kind} of contract:

\begin{description}[leftmargin=1.4em, style=nextline, nosep]
\item[Same-shape: \code{axpy}, \code{dot}.] One group $S$ shared across all tensor
positions. \lstinline[style=l4style]|axpy| returns shape $S$;
\lstinline[style=l4style]|dot| collapses \emph{every} axis to a \code{Scalar} (the
$S=[\,]$ case of a reduction). The contract is the simplest: the shapes must be
congruent, decided by \cref{def:b-resolve}.
\item[Contraction: \code{matvec}.] Two distinct axis variables $m,n$; the inner
$n$ appears on both the matrix's second axis and the vector's only axis, so they
must agree and are summed away; the outer $m$ survives as the result.
\code{matvec} is the rank-$1$/rank-$2$ instance of the general contraction whose
group-level form is the operator-apply of \cref{sec:b-operators}.
\item[Rank-raising: \code{outer}.] No axis is required to match; the result rank is
the \emph{sum} of the inputs'. The opposite move to \code{matvec}.
\item[Rank-moving: \code{reduce}, \code{broadcast}, \code{unfold}.]
\lstinline[style=l4style]|reduce| drops the labelled axis $k$ (\code{dot} is its
collapse-everything special case); \lstinline[style=l4style]|broadcast| stretches a
tensor to a larger \texttt{target} shape, with the matching side-condition
(\texttt{target} must be a broadcast extension of the source, decided by the
\lstinline[style=l4style]|DimExpr| solver); \lstinline[style=l4style]|unfold|
windows one axis into a new trailing \texttt{size} axis, the rank-raising mechanism
behind the matrix-free element gather.
\item[Structured kernel: \code{conv}.] The convolution contract is the most
structured: batch $N$ and input channels $C_i$ must agree between data and kernel,
the kernel's output-channel axis $C_o$ becomes the result's channel axis, and the
output spatial extent $S_o$ is a \lstinline[style=l4style]|DimExpr| function of the
input and kernel spatial extents $S_i$, $S_k$.
\end{description}

\section{Scalar promotion as subtyping}
\label{sec:b-promotion}

Everything so far concerned \emph{shapes}. One rule about \emph{element types}
threads through the contracts and deserves its own formal statement: scalar
promotion, the rule that lets a \emph{real} scalar join a \emph{complex} operation
without minting a second operator. Cast formally, it is a one-step \emph{subtyping}
relation \lstinline[style=l4style]|real <: complex|.

\begin{definition}[Scalar promotion subtyping]
\label{def:b-promotion}
The scalar element types carry the one-step subtype order
\[
  \texttt{real} \;\sqsubseteq\; \texttt{complex},
\]
witnessed by the coercion $\alpha \mapsto \alpha + 0i$ (attach a zero imaginary
part). The coercion is \emph{exact}: zero is representable exactly, so $\alpha$ and
$\alpha + 0i$ are the same number to the last bit, and every algebraic law of an
operator holds identically before and after promotion. The order is one-way: there
is no demotion \lstinline[style=l4style]|complex <: real|.
\end{definition}

As an inference rule, promotion is the standard \emph{subsumption} rule
specialized to the scalar lattice: a real-typed expression may be used where a
complex one is expected.

\[
  \dfrac{\Gamma \vdash e : \texttt{real}}
        {\Gamma \vdash e : \texttt{complex}}\quad(\textsc{sub-scalar})
  \qquad\qquad
  \dfrac{}{\texttt{real} \sqsubseteq \texttt{complex}}\quad(\textsc{sub-base})
\]

\noindent and it lifts to tensors elementwise: a
\lstinline[style=l4style]|Tensor[s]| over \texttt{real} subsumes to a
\lstinline[style=l4style]|Tensor[s]| over \texttt{complex} of the \emph{same}
shape---promotion changes only the element type, never the shape. The practical
payoff is the collapse of overloads: an operator like \code{axpy} need not come in
separate real-scalar and complex-scalar forms. When its tensors are \texttt{complex}
and a scalar argument arrives \texttt{real}, \textsc{sub-scalar} promotes the scalar
and the one complex \code{axpy} runs; its laws are written once.
\Cref{fig:b-lattice} draws the lattice.

\begin{figure}[t]
\centering
\begin{tikzpicture}[scale=1.0]
  \node[draw=simViolet, fill=simBgViolet, rounded corners=3pt, thick,
        minimum width=22mm, minimum height=8mm, font=\small] (c) at (0,1.5)
    {\texttt{complex}};
  \node[draw=simBlue, fill=simBgBlue, rounded corners=3pt, thick,
        minimum width=22mm, minimum height=8mm, font=\small] (r) at (0,0)
    {\texttt{real}};
  \draw[-{Stealth[length=2.6mm]}, simRust, thick] (r) -- (c)
    node[midway, right=2mm, font=\scriptsize\itshape, simRust, align=left]
    {\textsc{sub-scalar}:\\ $\alpha \mapsto \alpha + 0i$\\ (exact, no rounding)};
  \node[font=\scriptsize\itshape, simGray, align=right, text width=28mm]
    at (-2.4,0.75)
    {the subtype lattice\\ \lstinline[style=l4style]!real <: complex!};
  \draw[-{Stealth[length=2.2mm]}, simGray, dotted] (c.west) to[bend right=42]
    node[midway, left=1pt, font=\scriptsize, simGray] {\,no} (r.west);
\end{tikzpicture}
\caption[The scalar-promotion subtype lattice]{Scalar promotion is the one-step
subtype order \lstinline[style=l4style]!real <: complex!. Upward (rust): a
\texttt{real} scalar is \emph{promoted}---used where \texttt{complex} is
expected---by the exact coercion $\alpha\mapsto\alpha+0i$, the witness of
\textsc{sub-scalar}. The promotion attaches a zero imaginary part with no rounding,
so the operator's algebraic laws (\cref{app:laws}) are invariant. Downward (grey,
dotted): \emph{no} demotion---a \texttt{complex} scalar against a \texttt{real}
operation is a type error, because the imaginary part has nowhere to go. The
lattice runs one way, collapsing the real/complex operator overloads into one.}
\label{fig:b-lattice}
\end{figure}

\begin{pitfall}
Promotion lifts a \texttt{real} scalar \emph{into} a complex operation; it never
demotes, and it never touches a \emph{reduction's output} scalar. A
\lstinline[style=l4style]|dot :: Tensor[$S] -> Tensor[$S] -> Scalar| has no
\emph{incoming} scalar to promote---its scalar is the \emph{output}, fixed by the
operands' element type (complex operands give a complex result); likewise a norm's
result is real regardless. \textsc{sub-scalar} is about an incoming scalar
argument, not a result. If you find yourself wanting to ``promote'' a complex
coefficient down to a real result, the type is genuinely wrong: there is no
information-preserving demotion, and the calculus rejects it rather than silently
discard the imaginary part.
\end{pitfall}

\section{Tying back: where these rules are used}
\label{sec:b-tieback}

The rules of this appendix are not free-standing formalism; each is the precise
version of a discipline the main text relies on continuously.

\begin{itemize}[leftmargin=1.4em, nosep]
\item \textbf{To \cref{ch:tensors}.} The shape language (\cref{sec:b-shapelang})
and the primitive-contract table (\cref{fig:b-contracts}) formalize that chapter's
bracket notation and its picture-per-op treatment of \code{axpy}, \code{dot},
\code{matvec}, \code{outer}, \code{reduce}, \code{broadcast}. The \textsc{prim}
rule (\cref{sec:b-primrule}) is what its ``checked at every function boundary''
slogan means exactly.
\item \textbf{To \cref{ch:shapes}.} Named-group resolution (\cref{def:b-resolve})
is the formal version of that chapter's rank-agnostic congruence and its
boundary-gate figure; the range-first operator shapes (\cref{sec:b-operators}) and
the congruence-group-versus-fixed-axis distinction (\cref{sec:b-fixedaxes}) are its
\lstinline[style=l4style]|LinOp[(R: ...), (D: ...)]| and \texttt{E}/\texttt{L}/%
\texttt{P}/\texttt{C}/\texttt{G} material made formal; and scalar promotion
(\cref{sec:b-promotion}) is its \lstinline[style=l4style]!real <: complex! rule
cast as subtyping.
\item \textbf{To \cref{ch:matrixfree}.} The fixed-meaning element-local axes
(\cref{sec:b-fixedaxes}) are the shape vocabulary of the matrix-free kernel; the
square operator form \lstinline[style=l4style]|LinOp[(S: ...), $S]| is the shape a
Krylov solver demands of its operator, re-feedable at every power.
\item \textbf{To \cref{app:calculus} and \cref{app:laws}.} The typing judgment
$\Gamma \vdash e : \tau$ (\cref{sec:b-judgment}) shares its term grammar and
reduction rules with \cref{app:calculus}; the algebraic laws the \textsc{prim} rule
presupposes (and which scalar promotion leaves invariant) are collected in
\cref{app:laws}.
\end{itemize}

Read together, the appendix says one thing in four registers: a shape is a type, a
named group is a unification variable over axes of unknown rank, the
shape-contract rule \textsc{prim} is where shapes acquire force, and scalar
promotion is a one-step subtyping that keeps the algebra single-voiced across real
and complex.

\summarysection
This appendix supplies the formal object underneath \cref{ch:tensors,ch:shapes}. The
\emph{shape language} (\cref{sec:b-shapelang}) is a BNF grammar for a shape
expression $\sigma$: a bracketed tuple of items, each a literal extent, a named axis
variable, an algebraic combination ($+,-,\times,{}^\wedge$), the rank-wildcard
\lstinline[style=l4style]|...|, or a named group (binding
\lstinline[style=l4style]|(S: ...)|, use \lstinline[style=l4style]|$S|). Shapes are
compared modulo the \lstinline[style=l4style]|DimExpr| equational theory
(\cref{def:b-dimexpr})---commutative, associative $+$ and $\times$---so that
\lstinline[style=l4style]|h*w| matches \lstinline[style=l4style]|w*h|.
\emph{Named-group resolution} (\cref{def:b-resolve}) is a two-step unification:
\emph{bind} the group from the first labelled run, then at each use check rank
\emph{first} and extents modulo \lstinline[style=l4style]|DimExpr|---the precise
content of the congruence back-arrow, and the reason \lstinline[style=l4style]|Tensor[N]|
cannot match a rank-$2$ argument. Operator shapes name two groups range-first,
\lstinline[style=l4style]|LinOp[(R: ...), (D: ...)]|, collapsing to the square form
for endomorphisms; congruence groups are kept distinct from named axes of fixed
meaning (\texttt{E}/\texttt{L}/\texttt{P}/\texttt{C}/\texttt{G}). The
\emph{type-and-shape judgment} $\Gamma \vdash e : \tau$ (\cref{sec:b-judgment}) is
the standard typed-$\lambda$ form; shapes acquire force at exactly one rule, the
primitive shape-contract rule \textsc{prim}, which requires each argument's shape to
match the op's contract modulo the equational theory and group resolution, then
substitutes the resolved groups into the result. The \emph{primitive shape
contracts} (\cref{fig:b-contracts}) collect \code{axpy}, \code{dot}, \code{matvec},
\code{outer}, \code{reduce}, \code{unfold}, \code{broadcast}, and \code{conv} as a
reference table. Finally, \emph{scalar promotion}
(\lstinline[style=l4style]!real <: complex!, \cref{def:b-promotion}) is a one-step
subtyping witnessed by the exact coercion $\alpha\mapsto\alpha+0i$: a real scalar
joins a complex operation, never the reverse, and never on a reduction's output
scalar.

\furtherreading
The view of a matrix as a typed map between shaped spaces, rather than a grid of
numbers, is the organizing idea of Trefethen and Bau~\citep{trefethenbau1997}; their
treatment of matrix--vector products and the four fundamental subspaces is the
linear-algebra backdrop for the range-first operator-shape convention
(\cref{sec:b-operators}) and the contraction contract of \code{matvec}. The
inference-rule presentation of the typing judgment $\Gamma\vdash e:\tau$
(\cref{sec:b-judgment}), the subsumption rule behind scalar promotion
(\cref{sec:b-promotion}), and the general discipline of reading a type system as a
set of judgment rules follow the standard implementation-of-functional-languages
account of Peyton Jones~\citep{peytonjones2003}, whose treatment of type inference,
unification, and the typed-$\lambda$ core is the closest analogue to the
shape-as-type / congruence-as-unification reading developed here. The grammar and
reduction rules this appendix's typing judgment shares are collected in
\cref{app:calculus}; the algebraic laws the \textsc{prim} rule presupposes are
collected in \cref{app:laws}.
